% !TEX root = thesis.tex
\chapter{Conclusion and Future Work}
\label{sec:conclusion_and_future_work}

This chapter concludes the thesis by summarizing the developed online-learning architecture and its principal evaluation results. It then outlines a small number of extensions motivated directly by the limitations observed during implementation and hardware testing.

\section{Conclusion}

The objective of this thesis was to extend an \gls{snn} accelerator generated by Spiker+ with hardware support for \gls{sdsp}-based online learning. Achieving this objective required more than adding an isolated weight-update circuit. The generated neuron state had to be exposed to the learning logic, the originally fixed excitatory weights had to be moved into writable storage, and every completed update had to remain associated with the address from which its original weight was read. The resulting architecture retains the existing Spiker+ inference and synchronization structure while adding calcium state, local learning conditions, bounded weight arithmetic, and an address-aligned write-back path.

The accelerator was integrated as a reusable IP block in a PYNQ-Z1 system containing the required AXI control, DMA, and processor interfaces. Module-level simulation, network-level verification, and directed board tests were used to verify the individual learning functions and the complete data path. In particular, the teacher-control experiment showed that supervised operation caused a persistent and target-dependent change in later inference results. Reloading the bitstream restored the initial response, providing additional evidence that the observed change originated from the online synaptic state.

The final design meets timing at 100~MHz on the XC7Z020 and uses 5,056 LUTs, 5,294 flip-flops, and four BRAMs. The complete PYNQ application processed three continuous MNIST training epochs, corresponding to 180,000 image transactions, at an average end-to-end rate of 66.16 images/s. On the fixed class-balanced validation set, strict accuracy increased from 5.9\% before training to 39.6\% after the third epoch. The result confirms useful online adaptation, but it does not represent a converged or optimized classifier. The accuracy was still increasing, individual class responses remained unstable, and only one training run with one parameter configuration was evaluated.

The principal contribution of this work is therefore the design, integration, and board-level validation of persistent online weight adaptation within the Spiker+ hardware architecture. The experiments establish that the generated accelerator can be extended beyond inference-only operation and can retain learning results during continuous system execution. At the same time, the evaluation identifies clear opportunities for improving classification performance, experimental rigor, and system-level efficiency.

\section{Future Work}

Future work should build on the validated hardware path rather than changing several parts of the system simultaneously. Table~\ref{tab:future-work-directions} summarizes four directions and the measurements required to determine whether each extension is beneficial.

\begin{table}[htbp]
	\centering
	\caption{Prioritized Directions for Future Work}
	\label{tab:future-work-directions}
	\small
	\begin{tabularx}{\linewidth}{@{}p{0.20\linewidth}XX@{}}
		\toprule
		\textbf{Direction} & \textbf{Proposed extension} & \textbf{Evaluation target} \\
		\midrule
		Training and parameters & Longer training, controlled parameter sweeps, and repeated runs & Convergence, final test accuracy, and variation across runs \\
		Encoding and noise & Sparser or alternative encoding with controlled input and spike perturbations & Accuracy--activity trade-off, robustness, latency, and energy \\
		Network architecture & Add and train one hidden layer using a clearly defined local learning strategy & Accuracy improvement relative to memory, resource, and latency cost \\
		Measurement and runtime & Add cycle and weight observability, batch DMA transactions, and measure board power & Core latency, weight-level correctness, throughput, and energy per image \\
		\bottomrule
	\end{tabularx}
\end{table}


\paragraph{Longer Training and Parameter Evaluation.}
The first priority is to determine whether the observed learning trend continues and where it converges. Training should be extended substantially beyond five epochs, with checkpoints evaluated on a validation set that is separate from the final test set. The teacher strength, calcium thresholds and leakage, neuron threshold, rate-coding gain, input duration, and bistability mechanism should be varied systematically rather than together. Repeating the most promising configurations with different initial weights, image orders, and spike-generation seeds would make it possible to report variation as well as mean performance. This experiment would distinguish insufficient training from limitations imposed by the network and learning rule.

\paragraph{Input Encoding, Sparsity, and Noise Robustness.}
The current system uses stochastic rate coding for 100 time steps. Its computational activity is therefore determined partly by the number and distribution of generated input spikes. Rathi and Roy show that replacing Poisson-style input encoding with direct input encoding, together with neuron-parameter optimization, can substantially reduce latency and spike activity at comparable accuracy in directly trained deep SNNs \cite{Rathi2023DIETSNN}. Temporal coding provides another route to sparse activity, although the attainable accuracy and robustness depend on how information is represented in spike time. Park et al. demonstrate that neural coding choice affects the response of SNNs to spike deletion and timing jitter \cite{Park2021NoiseRobust}.

These results motivate, but do not guarantee, improvements to the local \gls{sdsp} architecture developed in this thesis. A suitable experiment should compare the present rate-encoding configuration with sparser rate settings and one alternative direct or temporal representation. For each configuration, the input-spike count, output activity, convergence behaviour, accuracy, processing time, and measured energy should be recorded. Clean inputs should additionally be compared with controlled pixel noise, spike deletion, and timing jitter. This would establish whether reducing redundant events improves efficiency without removing information required by online learning and whether input denoising produces a measurable improvement under the same training conditions.

\paragraph{Extension with a Hidden Layer.}
The single-layer 784--10 network was selected to keep the implementation compact and to provide a structural reference for comparison with TEDOP. Adding a hidden layer could improve the separation of digit features and reduce the dependence of classification on direct input--output associations. However, this extension would also increase synaptic storage, processing latency, and resource use. More importantly, a learning mechanism would be required for the hidden synapses because the class-dependent teacher signal currently acts at the output layer. Future work should therefore compare the present network with a 784--$H$--10 configuration under the same input and evaluation protocol, while explicitly defining whether the hidden layer is fixed, trained by local unsupervised plasticity, or coupled to an additional supervisory mechanism.

\paragraph{Measurement and Runtime Refinement.}
The current runtime is dominated by per-image Python control and contains a fixed software delay, while the reported power is a complete-device estimate dominated by the processing system. A hardware cycle counter and a debug-only weight read-back path would permit direct measurement of accelerator latency and individual synaptic changes. Reusing DMA buffers, batching several images per transfer, and replacing polling with interrupt-driven completion could reduce processor overhead. Finally, physical board-power measurements with controlled idle, inference, and learning phases would allow the influence of spike sparsity to be expressed as energy per image rather than inferred from a whole-device implementation estimate.

Together, these extensions would move the work from functional validation toward a more accurate and quantitatively optimized online-learning system. The immediate priority is a reproducible convergence and parameter study; architectural expansion and encoding changes should then be evaluated against that established baseline.
